\documentclass[11pt,a4paper]{article}

\usepackage[utf8]{inputenc}
\usepackage[T1]{fontenc}
\usepackage[spanish]{babel}

\usepackage{geometry}
\geometry{
    top=2.5cm,
    bottom=2.5cm,
    left=2.5cm,
    right=2.5cm
}

\usepackage{booktabs}
\usepackage{longtable}
\usepackage{array}
\usepackage{ragged2e}
\usepackage{enumitem}
\usepackage{hyperref}

\hypersetup{
    colorlinks=true,
    linkcolor=black,
    urlcolor=black,
    citecolor=black
}

\setlength{\parindent}{0pt}
\setlength{\parskip}{5pt}

\newcolumntype{P}[1]{>{\RaggedRight\arraybackslash}p{#1}}

\begin{document}

% =========================================================
% ENCABEZADO
% =========================================================

\begin{center}

{\Large\bfseries Informe de Correcciones}

\vspace{0.35cm}

{\itshape
``Evaluación Comparativa de Modelos de Series de Tiempo\\
con Estrategias de Ensamble para el Pronóstico\\
de la Demanda de Consumo Energético''
}

\vspace{0.25cm}

\textbf{Estudiante:} Juan David Piedrahita López

\textbf{Jurado:} Oveida Rosa Bustos Polo

\end{center}

\vspace{0.25cm}

\noindent\rule{\textwidth}{0.5pt}

\vspace{0.35cm}

Se presentan las correcciones realizadas a la monografía
a partir de las observaciones formuladas durante el proceso de evaluación.Las páginas indicadas corresponden a la versión corregida del documento en formato PDF.

\vspace{0.35cm}

% =========================================================
% TABLA DE CORRECCIONES
% =========================================================

\begin{longtable}{P{2.2cm} P{11.7cm}}

\toprule
\textbf{Página (PDF)} & \textbf{Corrección realizada} \\
\midrule
\endfirsthead



\toprule
\textbf{Página (PDF)} & \textbf{Corrección realizada} \\
\midrule
\endhead



\bottomrule
\endlastfoot

3 &
Se eliminaron puntos suspensivos adicionales. \\

\midrule

13 &
Se corrigió el enunciado relacionado con la mención de la metodología
Box--Jenkins. \\

\midrule

21 &
Se corrigió la expresión ``dado que la''. \\

\midrule

21 &
Se corrigió la sustitución utilizada en la expresión de la función de
autocorrelación, de acuerdo con la formulación presentada en Wei. \\

\midrule

23 &
Se corrigió el índice $\alpha$ posterior a la expresión 1.3. \\

\midrule

25 &
Se eliminó la sílaba ``es'' incluida por error. \\

\midrule

30 &
Se añadió la tilde faltante en la palabra ``parámetro''. \\

\midrule

32 &
\textbf{Observación 3.}
Se revisó la convención de signos utilizada para el polinomio de media móvil
estacional. En \textit{Time Series Analysis} de Wei, los coeficientes de dicho
polinomio aparecen con signo negativo en la expresión 8.33 (página 166),
mientras que en \textit{Time Series Analysis and Its Applications}, de
Shumway, se emplea una convención con signo positivo. \\

32 &
Se ajustó la formulación de la monografía para mantener consistencia con la
convención utilizada en esta última referencia. \\

\midrule

38 &
Se corrigió la palabra ``ensamble''. \\

\midrule

41 &
Se eliminó el artículo ``el'' repetido y se corrigió la palabra
``optimiza'' por ``optimizar''. \\

\midrule

44 &
Se corrigió la palabra ``algoritmo''. \\

\midrule

53 &
Se corrigió la expresión ``La expresión'' al inicio de la página. \\

\midrule

56 &
Se eliminó la palabra ``versión'' duplicada y se corrigió la palabra
``sigmoidal''. \\

\midrule

58 &
Se corrigió la expresión ``la matriz de los pesos''. \\

\midrule

64 &
Se corrigió la palabra ``inconveniente'' y se añadió el conector ``y'' en
la expresión ``\ldots\ a través del estado $h_t$ y del estado $h_{t-1}$''. \\

\midrule

65 &
Se corrigió la palabra ``trabajar''. \\

\midrule

70 &
Se corrigió la palabra ``mencionado''. \\

\midrule

71 &
Se corrigió la expresión ``donde $w_i$ es la \ldots''. \\

\midrule

73 &
Se corrigieron las expresiones ``AR-Net lineal'', ``Variables'', ``Forma''
y ``esencia''. \\

\midrule

83 &
\textbf{Observación 4.}
Se corrigió la fecha indicada en el pie de la Figura 6.1.
La fecha correcta corresponde al 1 de enero de 2015. \\

\midrule

85 &
Se corrigió la expresión ``adicional''. \\

\midrule

92--96 &
\textbf{Observación 2.}
Se corrigió la discrepancia identificada en las fechas utilizadas para la
evaluación del modelo XGBoost. La frontera temporal entre TrainG y TestF se
define inicialmente sobre la serie original mediante la partición temporal
90--10 utilizada para los demás modelos. \\

 &
Posteriormente, se construyen las variables temporales, los rezagos y los
promedios móviles. Las observaciones iniciales con valores \texttt{NaN}
generadas por estas transformaciones se eliminan sin modificar la fecha
previamente establecida para el inicio de TestF. \\

&
Como resultado, aunque el número de observaciones disponibles para el
entrenamiento de XGBoost se reduce debido a la construcción de variables
rezagadas, TestF conserva las mismas 58 observaciones y el mismo intervalo
temporal utilizado para evaluar los demás modelos. \\

 &
Se reejecutaron la optimización y el ajuste final de XGBoost bajo esta
partición. Los valores de algunas métricas e hiperparámetros presentan
variaciones respecto a la versión anterior; sin embargo, se conserva el
comportamiento general observado durante la validación temporal y XGBoost
continúa presentando el menor error entre los modelos individuales evaluados. \\

\midrule

93 &
Se eliminó el artículo ``a'' incluido por error. \\

\midrule

97 &
Se eliminó la expresión duplicada ``La partición contiene 517 observaciones
en TrainG y 58 en TestF''. \\

\midrule

109 &
\textbf{Observación 1.}
Se añadió en las conclusiones una aclaración sobre la limitación asociada al
cálculo de los pesos del esquema combinado utilizando el desempeño de los
modelos individuales sobre TestF. \\

 &
Debido a que este mismo conjunto se emplea posteriormente para evaluar la
predicción combinada, el resultado se interpreta como un análisis
retrospectivo y no como evidencia independiente de generalización sobre
observaciones futuras. \\

 &
Se señala además que una evaluación más rigurosa puede realizarse fijando
los pesos obtenidos previamente y evaluando posteriormente los modelos
individuales y el esquema combinado sobre un nuevo periodo temporal que no
haya intervenido en la definición de dichos pesos. \\

\midrule

-- &
Se corrigieron las métricas correspondientes al modelo NeuralProphet, dado
que los valores incluidos en la primera versión correspondían a una ejecución
anterior del notebook. \\

-- &
Se actualizaron las tablas, las métricas sobre TestF y la comparación entre
los modelos individuales. Con los valores corregidos, XGBoost presenta el
menor error entre los modelos individuales. \\

\end{longtable}

\end{document}